\input{preamble.tex}

\title{Лабораторная работа № 8. \\ Настройка сетевых сервисов. DHCP}
\author{Сущенко Алина Николаевна \\ НПИбд-01-23}

\institute{Российский университет дружбы народов имени Патриса Лумумбы}
\date{2026}

\begin{document}

\frame{\titlepage}

\begin{frame}
\frametitle{Цель работы}
\begin{itemize}
    \item Приобретение практических навыков по настройке динамического распределения IP-адресов посредством протокола DHCP в локальной сети.
\end{itemize}
\end{frame}

\begin{frame}
\frametitle{Логическая схема сети}
    \centering
    \includegraphics[width=0.9\textwidth]{../images/01.png}
    \captionof{figure}{Логическая схема локальной сети с DNS-сервером}
\end{frame}

\begin{frame}
\frametitle{Настройка DNS-сервера}
    \centering
    \includegraphics[width=0.9\textwidth]{../images/02.png}
    \captionof{figure}{Конфигурация сервера ansusenko-dns}
\end{frame}

\begin{frame}
\frametitle{Настройка DNS-сервера}
    \centering
    \includegraphics[width=0.9\textwidth]{../images/03.png}
    \captionof{figure}{Добавление ресурсных записей DNS}
\end{frame}

\frametitle{Проверка настроек DHCP-пулов}
    \centering
    \includegraphics[width=0.9\textwidth]{../images/04.png}
    \captionof{figure}{Настройка DHCP-пулов на маршрутизаторе}
\end{frame}

\begin{frame}
\frametitle{Просмотр информации о пулах DHCP}
    \centering
    \includegraphics[width=0.9\textwidth]{../images/05.png}
    \captionof{figure}{Результат выполнения команды \texttt{show ip dhcp pool}}
\end{frame}

\begin{frame}
\frametitle{Проверка привязок DHCP}
    \centering
    \includegraphics[width=0.9\textwidth]{../images/06.png}
    \captionof{figure}{Результат выполнения команды \texttt{show ip dhcp binding}}
\end{frame}

\begin{frame}
\frametitle{Настройка оконечных устройств на получение адресов по DHCP}
    \centering
    \includegraphics[width=0.9\textwidth]{../images/07.png}
    \captionof{figure}{Настройка интерфейса устройства \texttt{dk-ansusenko-01} на получение адреса по DHCP}
\end{frame}

\begin{frame}
\frametitle{Проверка доступности устройств}
    \centering
    \includegraphics[width=0.9\textwidth]{../images/08.png}
    \captionof{figure}{Результат проверки доступности между устройствами из разных подсетей}
\end{frame}

\begin{frame}
\frametitle{Проверка доступности устройств}
    \centering
    \includegraphics[width=0.9\textwidth]{../images/09.png}
    \captionof{figure}{Пинг устройств из разных подсетей}
\end{frame}

\begin{frame}
\frametitle{Режим симуляции DHCP}
    \centering
    \includegraphics[width=0.9\textwidth]{../images/10.png}
    \captionof{figure}{Движение пакетов в режиме симуляции при запросе адреса по DHCP}
\end{frame}


\begin{frame}
\frametitle{Контрольные вопросы}
\begin{block}{1. За что отвечает протокол DHCP?}
Протокол DHCP отвечает за автоматическую настройку сетевых параметров узлов в IP-сетях: IP-адрес, маску подсети, шлюз по умолчанию, DNS-сервера и другие параметры.
\end{block}
\begin{block}{2. Какие типы DHCP-сообщений передаются по сети?}
\begin{itemize}
    \item DHCPDISCOVER - запрос клиента на поиск DHCP-сервера
    \item DHCPOFFER - ответ сервера с предложением параметров
    \item DHCPREQUEST - подтверждение принятия предложения
    \item DHCPACK - подтверждение выделения адреса
    \item DHCPNAK - отклонение запроса
    \item DHCPRELEASE - освобождение адреса
\end{itemize}
\end{block}
\end{frame}

\begin{frame}
\frametitle{Контрольные вопросы}
\begin{block}{3. Какие параметры могут быть переданы в сообщениях DHCP?}
\begin{itemize}
    \item IP-адрес и маска подсети
    \item Адрес шлюза по умолчанию
    \item Адреса DNS-серверов
    \item Домен поиска DNS
    \item Время аренды адреса
    \item NTP-серверы
\end{itemize}
\end{block}
\begin{block}{4. Что такое DNS?}
DNS (Domain Name System) - распределённая иерархическая система для преобразования доменных имён в IP-адреса.
\end{block}
\end{frame}

\begin{frame}
\frametitle{Контрольные вопросы}
\begin{block}{5. Какие типы записи описания ресурсов есть в DNS и для чего они используются?}
\begin{itemize}
    \item \textbf{A} - связывает доменное имя с IPv4-адресом
    \item \textbf{AAAA} - связывает доменное имя с IPv6-адресом
    \item \textbf{CNAME} - создаёт псевдоним для другого имени
    \item \textbf{MX} - указывает почтовые серверы для домена
    \item \textbf{NS} - указывает DNS-серверы для зоны
    \item \textbf{PTR} - обратная запись, связывает IP-адрес с доменным именем
    \item \textbf{TXT} - хранит текстовую информацию
    \item \textbf{SOA} - содержит основную информацию о DNS-зоне
\end{itemize}
\end{block}
\end{frame}

\begin{frame}
\frametitle{Вывод}
\begin{itemize}
    \item В результате выполнения лабораторной работы приобретены практические навыки по настройке динамического распределения IP-адресов с использованием протокола DHCP.
    \item Настроен DNS-сервер \texttt{ansusenko-dns} с добавлением A-записей для домена \texttt{ansusenko.rudn.ru}.
    \item На маршрутизаторе \texttt{msk-ansusenko-gw-1} созданы DHCP-пулы для четырёх подсетей: \texttt{dk}, \texttt{departments}, \texttt{adm}, \texttt{other}.
    \item Оконечные устройства успешно получили IP-адреса по DHCP.
    \item Проверка доступности между устройствами из разных подсетей подтвердила корректность выполненных настроек.
\end{itemize}

\end{frame}

\end{document}
